% Source: Weinberg GR, physical PDF pp. 264--267
%         (printed pp. 241--244).
% Coverage: Section 9.8, Post-Newtonian Hydrodynamics; equations
%           (9.8.1)--(9.8.18).
% Figures/tables/footnotes: optional-reading footnote; reference marker 11;
%                           no figures or tables.
% Status: source-reviewed and compile-clean.
% Uncertainties: the final coefficient in the source's (9.8.13), repeated
%                on the left of (9.8.15), reads phi v^2. Dimensional
%                consistency and direct expansion of (9.8.1)--(9.8.3)
%                require rho v^2; this evident source typo is corrected.
%                The source also prints 1-v^2-2phi in (9.8.14), whereas
%                (9.8.11) and the stated sum of (9.3.2) and (9.3.4) require
%                1+v^2-2phi; that sign is corrected as well.
%
% MODERNIZATION CHECK: Divergences of momentum-flux dyadics are written
% with an explicit tensor product. Perturbative-order numerals are rendered
% as parenthesized superscripts.

\subsection{Post-Newtonian Hydrodynamics\optionalreading}\label{sec:9.8}

The post-Newtonian program outlined in Sections~\ref{sec:9.1}--\ref{sec:9.3}
would form an adequate basis for relativistic celestial mechanics if the
sun and planets could be regarded as point particles. However, this is not
the case; for instance, the tidal forces on the moon due to its finite size
are very much larger than the effects of the post-Newtonian corrections to
the earth's gravitational field. Often such finite-size effects may be
calculated to sufficient accuracy if we treat astronomical bodies as
composed of perfect fluids.\hyperref[ch9-ref-11]{\textsuperscript{11}} The
energy-momentum tensor is then given by Eq.~\eqref{eq:5.4.2}:
\begin{equation}
  T^{\mu\nu}
  =
  p g^{\mu\nu}+(p+\rho)U^\mu U^\nu,
  \tag{9.8.1}\label{eq:9.8.1}
\end{equation}
where \(p\) and \(\rho\) are the proper pressure and energy density, that
is, those measured by a locally comoving and freely falling observer, and
\(U^\mu\) is the four-vector velocity \(\dd x^\mu/\dd\tau\). (Of course,
\(p\) and \(\rho\) vanish except within the sun and planets.) To calculate
\(U^\mu\), we set
\begin{equation}
  \frac{U^i}{U^0}
  =
  \frac{\dd x^i}{\dd t}
  \equiv v^i,
  \tag{9.8.2}\label{eq:9.8.2}
\end{equation}
and calculate \(U^0\) from Eq.~\eqref{eq:9.2.2}:
\begin{equation}
  U^0
  =
  \frac{\dd t}{\dd\tau}
  =
  1-\phi+\frac12v^2+\order(\bar v^4).
  \tag{9.8.3}\label{eq:9.8.3}
\end{equation}

The program for calculating the motion of the fluid depends crucially on
whether there is an equation of state giving \(p\) as a function of
\(\rho\), as is the case for the cold degenerate fluids studied in
Chapter~11, or whether \(p\) depends on temperature as well. If the
pressure is a function of \(\rho\) alone, then our program is essentially
the same as in Section~\ref{sec:9.3}.

\medskip
\noindent\textbf{(A)} First solve the Newtonian problem. The pressure is to
be regarded as of order \(\bar v^2\bar M/\bar r^3\), so the necessary
components of the energy-momentum tensor are given by
Eqs.~\eqref{eq:9.8.1}--\eqref{eq:9.8.3} as
\begin{equation}
  T^{00(0)}=\rho,
  \tag{9.8.4}\label{eq:9.8.4}
\end{equation}
\begin{equation}
  T^{i0(1)}=\rho v^i,
  \tag{9.8.5}\label{eq:9.8.5}
\end{equation}
\begin{equation}
  T^{ij(2)}
  =
  p\delta_{ij}+\rho v^iv^j.
  \tag{9.8.6}\label{eq:9.8.6}
\end{equation}
The Newtonian equations of motion are provided by using
Eqs.~\eqref{eq:9.8.4}--\eqref{eq:9.8.6} in the mass- and
momentum-conservation equations \eqref{eq:9.3.2} and \eqref{eq:9.3.3}:
\begin{equation}
  \partial_t\rho
  +\boldsymbol\nabla\cdot(\rho\mathbf v)
  =0,
  \tag{9.8.7}\label{eq:9.8.7}
\end{equation}
\begin{equation}
  \partial_t(\rho\mathbf v)
  +\boldsymbol\nabla\cdot
   \bigl(\rho\,\mathbf v\otimes\mathbf v\bigr)
  =
  -\rho\boldsymbol\nabla\phi-\boldsymbol\nabla p,
  \tag{9.8.8}\label{eq:9.8.8}
\end{equation}
with \(p\) given as a function of \(\rho\) by the equation of state, and
with \(\phi\) determined by Poisson's equation \eqref{eq:9.3.12}:
\begin{equation}
  \nabla^2\phi=4\pi G\rho.
  \tag{9.8.9}\label{eq:9.8.9}
\end{equation}

\medskip
\noindent\textbf{(B)} Use the values of \(\rho,p,\mathbf v\), and \(\phi\)
determined in (A) to calculate the terms
\eqref{eq:9.8.4}--\eqref{eq:9.8.6} of \(T^{\mu\nu}\), and also to
calculate
\begin{equation}
  T^{00(2)}
  =
  \rho(v^2-2\phi).
  \tag{9.8.10}\label{eq:9.8.10}
\end{equation}

\medskip
\noindent\textbf{(C)} Use the results of (A) and (B) and
Eqs.~\eqref{eq:9.1.62} and \eqref{eq:9.1.65} to compute the
post-Newtonian fields \(\boldsymbol\zeta\) and \(\psi\).

\medskip
\noindent\textbf{(D)} Solve for \(\rho,p\), and \(\mathbf v\) in the
post-Newtonian approximation. The energy-momentum tensor is given to the
necessary order by Eqs.~\eqref{eq:9.8.1}--\eqref{eq:9.8.3} as
\begin{equation}
  T^{00(0)}+T^{00(2)}
  =
  \rho(1+v^2-2\phi),
  \tag{9.8.11}\label{eq:9.8.11}
\end{equation}
\begin{equation}
  T^{i0(1)}+T^{i0(3)}
  =
  (\rho+p+\rho v^2-2\phi\rho)v^i,
  \tag{9.8.12}\label{eq:9.8.12}
\end{equation}
\begin{equation}
  T^{ij(2)}+T^{ij(4)}
  =
  p\delta_{ij}(1+2\phi)
  +v^iv^j(p+\rho-2\phi\rho+\rho v^2),
  \tag{9.8.13}\label{eq:9.8.13}
\end{equation}
and the post-Newtonian equations of motion are obtained by using
Eqs.~\eqref{eq:9.8.11}--\eqref{eq:9.8.13} in the energy- and
momentum-conservation equations, \eqref{eq:9.3.2} plus
\eqref{eq:9.3.4}, and \eqref{eq:9.3.3} plus \eqref{eq:9.3.5}:
\begin{equation}
  % VERIFY SOURCE: the printed -v^2 conflicts with (9.8.11) and with the
  % explicit sum of (9.3.2) and (9.3.4); the required sign is +v^2.
  \begin{split}
  &\partial_t\!\left[\rho(1+v^2-2\phi)\right]\\
  &\quad
  +\boldsymbol\nabla\cdot
  \left[\mathbf v(\rho+p+\rho v^2-2\phi\rho)\right]
  =
  \rho\,\partial_t\phi,
  \end{split}
  \tag{9.8.14}\label{eq:9.8.14}
\end{equation}
\begin{equation}
  \begin{split}
  &\partial_t
  \left[\mathbf v(\rho+p+\rho v^2-2\phi\rho)\right]\\
  &\quad
  +\boldsymbol\nabla\cdot
  \left[
    (\mathbf v\otimes\mathbf v)
    (p+\rho-2\phi\rho+\rho v^2)
  \right]\\
  &=
  -\boldsymbol\nabla[p(1+2\phi)]
  -\rho\boldsymbol\nabla(\phi+2\phi^2+\psi)
  -\rho\,\partial_t\boldsymbol\zeta\\
  &\quad
  -\rho(v^2-2\phi)\boldsymbol\nabla\phi
  +\rho\mathbf v\times
   (\boldsymbol\nabla\times\boldsymbol\zeta)
  +4\rho\mathbf v\,\partial_t\phi\\
  &\quad
  -(3p+\rho v^2)\boldsymbol\nabla\phi
  +4\phi\rho\boldsymbol\nabla\phi
  +4\rho\mathbf v
   (\mathbf v\cdot\boldsymbol\nabla\phi).
  \end{split}
  \tag{9.8.15}\label{eq:9.8.15}
\end{equation}

\medskip
\noindent\textbf{(E)} And so on.

Matters are more complicated when the temperature is an independent
variable. We now need one additional equation at every stage in the
calculation, which is provided for us by an equation of continuity:
\begin{equation}
  \partial_\mu\left(\sqrt{-g}\,\mu U^\mu\right)=0,
  \tag{9.8.16}\label{eq:9.8.16}
\end{equation}
where \(\mu\) is a rest-mass density proportional to the number density of
particles in the fluid. [Compare Eq.~\eqref{eq:5.2.14}.] It may be assumed
that the pressure is given by the equation of state as a function of both
\(\mu\) and an energy density \(\epsilon=\order(\bar v^2)\), defined by
\begin{equation}
  T^{00}\equiv\mu U^0+\epsilon.
  \tag{9.8.17}\label{eq:9.8.17}
\end{equation}
Our equations are then the continuity equation \eqref{eq:9.8.16}, the
momentum-conservation equation \(\nabla_\mu T^{\mu i}=0\), and an
energy-conservation equation, which after subtracting Eq.~\eqref{eq:9.8.16}
may be written
\begin{equation}
  \partial_t\left(\sqrt{-g}\,\epsilon\right)
  +\partial_i
   \left\{\sqrt{-g}\,[T^{i0}-\mu U^i]\right\}
  =
  -\sqrt{-g}\,\Gamma^0{}_{\mu\nu}T^{\mu\nu}.
  \tag{9.8.18}\label{eq:9.8.18}
\end{equation}

However, now we use the energy-conservation equation to one higher order
in \(\bar v^2\) at every step. In the Newtonian approximation we use the
continuity equation to order \(\bar v\), the momentum-conservation equation
to order \(\bar v^2\), and the energy-conservation equation to order
\(\bar v^3\), while in the post-Newtonian approximation we use continuity
to order \(\bar v^3\), momentum conservation to order \(\bar v^4\), and
energy conservation to order \(\bar v^5\). Without writing out these
equations in detail, we may note that this program is possible, because in
the Newtonian calculation we need \(\Gamma^{0(3)}{}_{00}\) and
\(\Gamma^{0(2)}{}_{i0}\), which are given by Eqs.~\eqref{eq:9.1.71} and
\eqref{eq:9.1.72} in terms of \(\phi\) alone, while in the post-Newtonian
calculation we also need \(\Gamma^{0(5)}{}_{00}\),
\(\Gamma^{0(4)}{}_{i0}\), and \(\Gamma^{0(3)}{}_{ij}\), which are given
by Eqs.~\eqref{eq:9.1.73}--\eqref{eq:9.1.75} in terms of the
post-Newtonian fields.
